%% main-llmxive.tex — PROJ-562's paper, re-styled in the llmXive house style.
%%
%% This is the "publication-ready" version of the paper that llmXive shows on
%% its dashboard. It wraps the original main.tex body in the llmxive document
%% class — typography (Fraunces + JetBrains Mono) and color (Dartmouth green)
%% match the website. The original main.tex (and main.pdf, if you compile it
%% with pdflatex) is kept for reference + provenance.
%%
%% Compile from the llmXive repo root:
%%   cd /path/to/llmXive
%%   xelatex -output-directory=projects/PROJ-562-…/paper/pdf \
%%           projects/PROJ-562-…/paper/source/main-llmxive.tex

\documentclass[11pt]{llmxive}

% Re-declare the original's figure-number aliases (the body uses them in
% \cref / \autoref calls — they're cross-reference labels, not figure data).
\newcommand{\crossentropyContent}{1}
\newcommand{\crossentropyFunction}{2}
\newcommand{\crossentropyPOS}{3}
\newcommand{\ttestsContent}{4}
\newcommand{\ttestsFunction}{5}
\newcommand{\ttestsPOS}{6}
\newcommand{\confusion}{7}
\newcommand{\mds}{8}
\newcommand{\authortableContent}{1}
\newcommand{\authortableFunction}{2}
\newcommand{\authortablePOS}{3}

% The original is double-spaced; this re-style is single-spaced (more
% web-aligned). Comment-in to restore double-spacing:
% \usepackage{setspace}\doublespacing

\title{A Stylometric Application of Large Language Models}
\author{Harrison F. Stropkay, Jiayi Chen, Mohammad J. Latifi,\\
        Daniel N. Rockmore, and Jeremy R. Manning\\[0.3em]
        \normalsize\itshape Dartmouth College, Hanover, NH 03755, USA}
\date{October 2025}
\arxivid{2510.21958}
\githubrepo{https://github.com/ContextLab/llm-stylometry}
\journal{llmXive --- automated scientific discovery}

\begin{document}
\maketitle

\begin{abstract}
We show that large language models (LLMs) can be used to distinguish the
writings of different authors. Specifically, an individual GPT-2 model,
trained from scratch on the works of one author, will predict held-out text
from that author more accurately than held-out text from other authors. We
suggest that, in this way, a model trained on one author's works embodies the
unique writing style of that author. We first demonstrate our approach on
books written by eight different (known) authors. We also use this approach
to confirm R.~P.~Thompson's authorship of the well-studied
15\textsuperscript{th} book of the \textit{Oz} series, originally attributed
to F.~L.~Baum.
\end{abstract}

\section{Introduction}

Herein we introduce \textit{predictive comparison}, a new LLM-based relative
stylometric measure. It derives from a simple idea, that if an LLM can be
trained to write like---i.e., in the style of---a given author by training on
their work~\citep[e.g., ][]{Mikr25}, then the degree to which such a model
can be used to predict held-out text from any \emph{candidate} author should
reflect the relative similarity of the styles of the two authors.

For ease of exposition, throughout this paper we focus on prose written by
known novelists. To this end, we adapted and trained a series of small GPT-2
models~\citep{RadEtal19} to write in the style of any of $N=8$ different
authors. Specifically, for each author, we trained one model on a training
partition of that author's complete works, and evaluated each model's
predictive accuracy on held-out passages from each of the eight authors.

\section*{Acknowledgments}

The authors thank Wenjia Tracy Yu and Lucy L.~W.~Owen for useful comments and
discussion. We acknowledge support from NSF grant number CAREER-2145172.

\section*{Code availability}

Code: \url{https://github.com/ContextLab/llm-stylometry}. See the
project's README for setup instructions.

\section*{Note}

\textbf{This is an abbreviated re-styled excerpt} of the full paper, produced
as a demonstration of the llmXive house style. The full content
(Methods, Results, Discussion, all figures, and the complete bibliography) is
available in the original on arXiv at
\href{https://arxiv.org/abs/2510.21958}{arXiv:2510.21958}.

\end{document}
